% from ega4-2-fr.tex
\oldpage[IV-2]{178}
(i) Si $f$ est radiciel au point $y$ et $g$ radiciel au point $z$,
$g^{-1}(z)$ est réduit à $y$ et
$f^{-1}(g^{-1}(z))=f^{-1}(y)$ vide ou réduit à un seul point $x$; en
outre $k(x)$ est extension radicielle de $k(y)$ et $k(y)$ extension
radicielle de $k(z)$, donc $k(x)$ est extension radicielle de $k(z)$.
Inversement, supposons $f$ surjectif; si $g\circ f$ est radiciel au point
$z=g(y)$, $g^{-1}(z)$ se réduit au seul point $y$, sans quoi
$f^{-1}(g^{-1}(z))$ aurait au moins deux points distincts; en outre,
$f^{-1}(y)=f^{-1}(g^{-1}(z))$ se réduit à un seul point $x$, et par
hypothèse on a $k(z)\subset k(y)\subset k(x)$, et $k(x)$ est radiciel sur
$k(z)$, donc $k(y)$ est radiciel sur $k(z)$ et $k(x)$ radiciel sur $k(y)$.

(ii) Soient $g:f^{-1}(y)\to\operatorname{Spec}(k(y))$,
$g':f'^{-1}(y')\to\operatorname{Spec}(k(y'))$ les morphismes images
réciproques de $f$ et $f'$ respectivement; comme
$f'^{-1}(y')=f^{-1}(y)\otimes_{k(y)}k(y')$ (I, 3.6.4), $g'$ est image
réciproque de $g$ et l'assertion se réduit à (2.6.1, (v)).
\end{proof}

Soit alors $X$ un préschéma réduit n'ayant qu'un nombre fini de
composantes irréductibles, et soit $X'$ son normalisé (II, 6.3.8); on sait
(loc. cit.) que le morphisme canonique $f:X'\to X$ est surjectif. La
définition (6.15.1) montre alors que pour que $X$ soit géométriquement
unibranche en un point $x$, il faut et il suffit que $f$ soit radiciel en ce
point; pour que $X$ soit géométriquement unibranche, il faut et il suffit
donc que $f$ soit radiciel.

\begin{env}[6.15.4]
\EGAReaderTarget{ega4.c04.statement.6-15-4}{IV.2.s6.15.4-fr}
Généralisant la définition donnée dans (I, 2.2.9), nous dirons que, pour
deux préschémas réduits $X$, $Y$, un morphisme $f:X\to Y$ est
\emph{birationnel} si la restriction de $f$ à l'ensemble des points
maximaux de $X$ est une bijection de cet ensemble sur l'ensemble des points
maximaux de $Y$, et si, pour tout point maximal $y$ de $Y$, le morphisme
$X\times_Y\operatorname{Spec}(k(y))\to\operatorname{Spec}(k(y))$ déduit de
$f$ est un isomorphisme (autrement dit, si la fibre $f^{-1}(y)$ est réduite
à un seul point $x$ (maximal dans $X$) et si l'homomorphisme
$k(y)\to k(x)$ déduit de $f$ est bijectif); cette notion se réduit à celle
de (I, 2.2.9) lorsque $X$ et $Y$ n'ont qu'un nombre fini de composantes
irréductibles.
\end{env}

\begin{lemma}[6.15.4.1]
\label{IV.2.s6.15.4.1-fr}
Soit $f:X\to Y$ un morphisme, $g:Y'\to Y$ un morphisme plat,
$X'=X\times_Y Y'$, $f'=f_{(Y')}:X'\to Y'$. Si $f_{\mathrm{red}}$ est
birationnel, il en est de même de $f'_{\mathrm{red}}$.
\end{lemma}

\begin{proof}
En effet, si $y'$ est un point maximal de $Y'$, on sait (2.3.4, (ii)) que
$y=g(y')$ est un point maximal de $Y$; il existe un seul point $x$ de $X$
dans $f^{-1}(y)$ par hypothèse; de plus $x$ est maximal dans $X$ et tel que
$k(x)=k(y)$. Il résulte donc de (I, 3.4.9) que $f'^{-1}(y')$ est réduit à
un seul point $x'$ et que $k(x')=k(y')$. On conclut le raisonnement en
remarquant que d'après (2.3.7, (ii)), toute composante irréductible de $X'$
domine une composante irréductible de $Y'$.
\end{proof}

\begin{proposition}[6.15.5]
\label{IV.2.s6.15.5-fr}
Soit $f:X\to Y$ un morphisme tel que $f_{\mathrm{red}}$ soit entier et
birationnel (6.15.4), donc surjectif.
\begin{enumerate}
\item[(i)] Pour que $Y$ soit géométriquement unibranche en un point $y$,
il faut et il suffit que $f$ soit radiciel au point $y$ et que $X$ soit
géométriquement unibranche en l'unique point $x$ de $f^{-1}(y)$.
\item[(ii)] Pour que $Y$ soit géométriquement unibranche, il faut et il
suffit que $X$ soit géométriquement unibranche et que $f$ soit radiciel.
\end{enumerate}
\end{proposition}

\begin{proof}
On peut évidemment se borner au cas où $X$ et $Y$ sont réduits, le fait
que $f$ soit surjectif résulte de ce que $f$ est un morphisme fermé (II,
6.1.10) et que $f(X)$ contient par hypothèse tous les points maximaux de
$Y$. Cela montre aussitôt que (i) entraîne (ii). Utilisant (I, 3.6.5) et
(II, 6.1.5), on peut supposer, pour prouver (i), que
$Y=\operatorname{Spec}(\mathcal O_{Y,y})$, autrement dit
$Y=\operatorname{Spec}(A)$, où $A$ est local. Comme $f$ est affine, on a
aussi $X=\operatorname{Spec}(B)$. Si $A$ est géométriquement unibranche,
il est intègre, donc il en est de même de $B$ puisque $f$ est birationnel;
inversement, si $f$ est radiciel au point $y$ et $X$ géométriquement
unibranche au point $x$, $B$ n'a qu'un seul idéal maximal (puisqu'il est
entier sur $A$ et que tout idéal maximal de $B$ est donc au-dessus de celui
de $A$), autrement dit $B$ est un anneau local, et dire que $X$ est
géométriquement unibranche au point $x$ signifie que $B$ est
géométriquement unibranche, donc intègre. Comme $f$ est dominant par
hypothèse et que $A$ est réduit, on en conclut que $A\subset B$ (I,
1.2.7), donc $A$ est aussi intègre. Ainsi, pour prouver (i), on peut se
borner au cas où $A$ et $B$ sont locaux, intègres, $A$ étant contenu dans
$B$ et ayant même corps des fractions, et par suite (puisque $B$ est entier
sur $A$) même clôture intégrale $C$. Mais alors l'assertion résulte de
(6.15.3.1, (i)) appliqué aux morphismes
$\operatorname{Spec}(C)\to\operatorname{Spec}(B)$ et
$\operatorname{Spec}(B)\to\operatorname{Spec}(A)$, vu l'interprétation
donnée dans (6.15.3) de la propriété d'être géométriquement unibranche en
un point.
\end{proof}

